% ============================================================
%  GBM Multimodal Clustering — 3-Slide Technical Presentation
%  Compile with: pdflatex presentation_slides.tex
% ============================================================
\documentclass[aspectratio=169, 10pt]{beamer}

% ---- Theme & Colors ----
\usetheme{Madrid}
\usecolortheme{whale}
\definecolor{DarkTeal}{RGB}{0, 77, 96}
\definecolor{Highlight}{RGB}{204, 60, 60}
\definecolor{GoodGreen}{RGB}{34, 139, 34}
\definecolor{WarnOrange}{RGB}{210, 120, 20}
\setbeamercolor{frametitle}{bg=DarkTeal, fg=white}
\setbeamercolor{title}{bg=DarkTeal, fg=white}
\setbeamercolor{block title}{bg=DarkTeal, fg=white}
\setbeamercolor{item}{fg=DarkTeal}
\setbeamertemplate{navigation symbols}{}
\setbeamertemplate{footline}[frame number]

\usepackage{booktabs}
\usepackage{graphicx}
\usepackage{xcolor}
\usepackage{enumitem}
\setlist[itemize]{nosep, leftmargin=*, topsep=2pt, itemsep=1pt, parsep=0pt}

\title[GBM Clustering]{\large Multi-Modal Clustering Identifies a High-Risk GBM Subtype}
\subtitle{Combining Brain MRI Features with Clinical Data}
\author{GBM Research Team}
\date{}

\begin{document}

% ============================================================
% SLIDE 1: Project Overview & Problem Statement
% ============================================================
\begin{frame}{\textbf{Project Overview \& Problem Statement}}
\small

\textbf{The Problem in plain language:}
\begin{itemize}
    \item Glioblastoma (GBM) is the \textbf{most aggressive brain cancer} with average survival at roughly 15 months
    \item Patients are treated similarly, but not all tumours behave the same way so identifying the \textbf{worst subtype} early could help give personalised treatment.
\end{itemize}

\vspace{0.15cm}
\textbf{What data do we have?}
\begin{itemize}
    \item \textbf{Brain MRI scans} : tumour sub region ratios (necrosis, edema, enhancing core) and dominant brain lobe
    \item \textbf{Clinical records} : age, genetic markers (MGMT methylation, IDH mutation), survival outcomes
    \item \textbf{Two independent hospital cohorts:} UCSF (501 patients) and UPenn (611 patients)
\end{itemize}

\vspace{0.15cm}
\textbf{Objective:}
\begin{center}
\fcolorbox{DarkTeal}{white}{\parbox{0.88\textwidth}{\centering\small
\textit{Use unsupervised machine learning to discover patient subgroups from combined MRI and clinical features, then validate whether these subgroups have genuinely different survival outcomes.}
}}
\end{center}

\end{frame}

% ============================================================
% SLIDE 2: Methodology
% ============================================================
\begin{frame}{\textbf{Methodology}}
\footnotesize

\begin{columns}[T]
\begin{column}{0.48\textwidth}

\textbf{Step 1: Feature Extraction \& Cleaning}
\begin{itemize}
    \item Extract 31 features per patient from MRI and clinical records
    \item It Includes: tumour sub region ratios, dominant lobe, tumour size, age, genetic markers
\end{itemize}
\vspace{0.1cm}

\textbf{Step 2: Anti-Leakage Split (70/30)}
\begin{itemize}
    \item Stratified train test split on MGMT status and survival
    \item Imputation \& scaling learned \textbf{only from training data}
    \item \textit{prevents data leakage}
\end{itemize}
\vspace{0.1cm}

\textbf{Step 3: Spectral Clustering}
\begin{itemize}
    \item Group patients by feature similarity; test $k$ (cluster size) for 2-6 clusters
    \item \textit{Goal is to find natural groupings where patients with similar characteristics are grouped together}
\end{itemize}

\end{column}
\begin{column}{0.48\textwidth}

\textbf{Step 4: Nested Model Selection}
\begin{itemize}
    \item Best $k$ chosen on an \textbf{inner split} of the training set (unused test set)
    \item 6 weighted criteria: cluster quality, stability, survival separation, biological coherence
\end{itemize}
\vspace{0.1cm}

\textbf{Step 5: Validation}
\begin{itemize}
    \item Assign held out test patients to nearest cluster centre
    \item Check if survival differences \textbf{replicate} on unseen data
\end{itemize}
\vspace{0.1cm}

\textbf{Step 6: Repeated Testing}
\begin{itemize}
    \item \textbf{Repeated Validation:} re run across 30 random splits
    \item \textbf{Bootstrap Stability:} re cluster 100 times to confirm reproducibility
\end{itemize}

\end{column}
\end{columns}

\end{frame}

% ============================================================
% SLIDE 3: Results & Interpretation
% ============================================================
\begin{frame}{\textbf{Results \& Interpretation}}
\footnotesize

\begin{columns}[T]
\begin{column}{0.47\textwidth}

\textbf{UCSF — Strong Validation} \textcolor{GoodGreen}{\textbf{$\checkmark$}}

\vspace{0.08cm}
{\small
\begin{tabular}{@{}l r@{}}
\toprule
\textbf{Metric} & \textbf{Value} \\
\midrule
Selected $k$ (Cluster Size)       & 2 \\
Train log-rank $p$                & 4.79\,$\times$\,$10^{-11}$ \\
\textbf{Test log-rank $p$}        & \textcolor{GoodGreen}{\textbf{1.70\,$\times$\,$10^{-6}$}} \\
Lobe match on test                & \textcolor{GoodGreen}{Yes (frontal)} \\
NC/EN rank on test                & \textcolor{GoodGreen}{\#1 of 2} \\
High-risk: train/test size        & 83\% / 76\%  \\
\bottomrule
\end{tabular}
}

\end{column}
\begin{column}{0.47\textwidth}

\textbf{UPenn — Partial Replication} \textcolor{WarnOrange}{\textbf{$\sim$}}

\vspace{0.08cm}
{\small
\begin{tabular}{@{}l r@{}}
\toprule
\textbf{Metric} & \textbf{Value} \\
\midrule
Selected $k$ (Cluster Size)        & 4 \\
Train log-rank $p$                 & 0.023 \\
\textbf{Test log-rank $p$}         & \textcolor{Highlight}{0.20 (n.s.)} \\
Lobe match on test                 & \textcolor{Highlight}{No (parietal)} \\
NC/EN rank on test                 & \#2 of 4 \\
High-risk: train/test size         & 16\% / 14\% \\
\bottomrule
\end{tabular}
}

\end{column}
\end{columns}

\vspace{0.1cm}
\begin{center}
\textbf{Validation (30 random splits):}
\end{center}
\vspace{-0.4cm}

\begin{columns}[T]
\begin{column}{0.47\textwidth}
\begin{itemize}
    \item Test survival significant: \textcolor{GoodGreen}{\textbf{30/30}} runs
    \item Cluster properties stable: \textcolor{GoodGreen}{29/30} runs
    \item Bootstrap consensus ARI\,=\,\textbf{0.94} {\scriptsize(1.0 = perfect)}
\end{itemize}
\end{column}
\begin{column}{0.47\textwidth}  
\begin{itemize}
    \item Test survival significant: \textcolor{WarnOrange}{\textbf{7/30}} runs
    \item Temporal phenotype replicated: \textcolor{GoodGreen}{\textbf{22/27}} runs
    \item Bootstrap consensus ARI\,=\,\textbf{0.80}
\end{itemize}
\end{column}
\end{columns}

\vspace{0.1cm}
\rule{\textwidth}{0.3pt}
{\footnotesize
\textbf{Log-rank $p$-value:} tests if clusters have different survival curves; $p < 0.05$ = significant.\;
\textbf{NC/EN rank:} shows if the high-risk group have the highest necrosis ratio on unseen data\;
\textbf{ARI:} measures clustering reproducibility; 1.0 = identical, 0 = random.
}

\vspace{0.08cm}
\fcolorbox{DarkTeal}{white}{\parbox{0.96\textwidth}{\centering\footnotesize
\textbf{Takeaway:} UCSF strongly identifies a high necrosis, poor survival GBM subtype validated across all 30 test sets. UPenn replicates the tumour phenotype but not survival separation
}}

\end{frame}

\end{document}